\section*{Solution}

% In what follows we give a complete and self-contained proof of the theorem.
% The argument follows the strategy laid out in the problem: we develop the
% cut/flow lattice machinery, prove a projection length formula, deduce the
% irreducibility of certain dual vectors via apex trees, and combine these with
% a structural dichotomy for reducible vertices.

\begin{thm}\label{thm:main}
Let $T$ be a finite weighted tree with vertex set $V$, edge set $E$, and
weight function $w:V\to\bZ$. Let $L(T)$ be the free abelian group on $V$ with
the symmetric bilinear form determined by
\[
u\cdot v=\begin{cases} w(u),& u=v,\\ -1,& (u,v)\in E,\\ 0,& u\neq v,\ (u,v)\notin E.\end{cases}
\]
Suppose the form is positive definite and that there exists a single vertex $v$
with $w(v)<d(v)$ (where $d(v)$ is the degree of $v$). Then $L(T)$ contains a
vertex element $u\in V$ that is irreducible, i.e.\ there do not exist nonzero
$a,b\in L(T)$ with $u=a+b$ and $a\cdot b\geq 0$.
\end{thm}

% Throughout, $|x|^2 := x\cdot x$.  A vertex $v$ is *bad* if $w(v)<d(v)$.

\subsection*{Cut and flow lattices}

% We recall the standard definitions of cut and flow lattices on a finite graph,
% and record the facts we need.

Let $G$ be a finite loopless graph (parallel edges allowed) with vertex set
$V(G)$ and edge set $E(G)$. Write $e(v,w)$ for the number of edges joining
distinct vertices $v,w$. Fix an orientation of the edges; this makes $G$ a
$1$-dimensional CW complex with cellular chain complex
\[
0\longrightarrow C_1(G)\xrightarrow{\ \partial\ } C_0(G)\longrightarrow 0,
\qquad \partial(e)=v-w \text{ for } e \text{ oriented from } w \text{ to } v.
\]
We endow $C_1(G)$ and $C_0(G)$ with the inner products in which $E(G)$ and
$V(G)$ are orthonormal bases. The adjoint $\partial^{*}:C_0(G)\to C_1(G)$ is
$\partial^{*}v=\sum_{e}(\partial e\cdot v)e$, and one computes
\begin{equation}\label{eq:cob}
\partial^{*}u\cdot \partial^{*}v=\begin{cases} d(u),& u=v,\\ -e(u,v),& u\neq v.\end{cases}
\end{equation}
The *cut lattice* is $\Cut(G)=\im\partial^{*}\subseteq C_1(G)$ and the *flow
lattice* is $\Flow(G)=\ker\partial\subseteq C_1(G)$. Since $\partial^{*}$ is
adjoint to $\partial$, $\Cut(G)$ and $\Flow(G)$ are orthogonal complements of
each other inside the (rational) vector space $C_1(G)\otimes\bQ$.

% Both are primitive sublattices of the self-dual lattice $C_1(G); we record this.

\begin{lem}\label{lem:primitive}
$\Cut(G)$ and $\Flow(G)$ are primitive sublattices of $C_1(G)$, and they are
mutually orthogonal complements. Consequently the orthogonal direct sum
$\Cut(G)\oplus\Flow(G)\hookrightarrow C_1(G)$ induces, after dualizing and using
the self-duality $C_1(G)\cong C_1(G)^{*}$, an inclusion
$C_1(G)\hookrightarrow \Cut(G)^{*}\oplus\Flow(G)^{*}$. Every $x\in C_1(G)$ thus
has a unique expression $x=x_C+x_F$ with $x_C\in\Cut(G)^{*}$,
$x_F\in\Flow(G)^{*}$, where $x_C$ (resp.\ $x_F$) is the orthogonal projection of
$x$ onto $\Cut(G)\otimes\bQ$ (resp.\ $\Flow(G)\otimes\bQ$).
\end{lem}

\begin{proof}
Orthogonality is the adjointness statement above. For primitivity of
$\Flow(G)=\ker\partial$: if $x\in C_1(G)$ and $nx\in\Flow(G)$ for some integer
$n\neq0$, then $n\,\partial x=0$ in the torsion-free group $C_0(G)$, so
$\partial x=0$ and $x\in\Flow(G)$. Thus
$C_1(G)\cap(\Flow(G)\otimes\bQ)=\Flow(G)$, which is one of the equivalent
definitions of primitivity. For $\Cut(G)=\im\partial^{*}$: the cokernel
$C_1(G)/\Cut(G)$ is torsion-free because $\Cut(G)$ is the orthogonal complement
of the primitive sublattice $\Flow(G)$ (indeed
$C_1(G)\cap(\Cut(G)\otimes\bQ)=\Cut(G)$ for the same reason: if $nx$ lies in the
orthogonal complement of $\Flow(G)$ then so does $x$). Hence $\Cut(G)$ is
primitive as well.

For the last assertion, recall that for a primitive sublattice $P\subseteq M$ of
a self-dual lattice $M$ with orthogonal complement $Q$, the restriction map
$M=M^{*}\to P^{*}$ and $M\to Q^{*}$ are surjective; combining them gives the
injection $M\hookrightarrow P^{*}\oplus Q^{*}$ realized by orthogonal
projection. The orthogonal projection of $x$ onto $P\otimes\bQ$ is precisely the
$P^{*}$-component, because for $p\in P$ we have $x\cdot p=x_C\cdot p$. This
proves the claim.
\end{proof}

% Note for later use: since $C_1(G) \to \Cut(G)^*$ (orthogonal projection) is
% surjective, every element of \Cut(G)^* arises as x_C for some x in C_1(G).

\subsection*{The projection length formula}

% We now prove the key length formula via a "completing the square" argument.

Fix a connected graph $G$ with oriented edges. For an ordered partition
$(A,B)$ of $V(G)$, let $E(A,B)$ be the set of edges with one endpoint in each
part, and set
\[
\chi_{A,B}=\sum_{e\in E(A,B)}\varepsilon_{A,B,e}\,e,\qquad
\varepsilon_{A,B,e}=\begin{cases}+1,& e \text{ oriented from } A \text{ to } B,\\ -1,& e \text{ oriented from } B \text{ to } A.\end{cases}
\]
Then $\chi_{B,A}=-\chi_{A,B}$, and
\begin{equation}\label{eq:chi}
\chi_{A,B}=-\sum_{v\in A}\partial^{*}v\ \in\ \Cut(G),
\end{equation}
% indeed each edge inside A or inside B cancels, and each edge of E(A,B)
% contributes the sign \varepsilon_{A,B,e}.
which one checks edge by edge.

A *two-component spanning forest* of $G$ is an unordered pair $\{A,B\}$ where
$A,B$ are vertex-disjoint subtrees whose vertex sets partition $V(G)$. For
$x=\sum_e x_e e\in C_1(G;\bR)$ define its *square-flux*
\[
\cF_{\{A,B\}}(x)=|\chi_{A,B}\cdot x|^2
=\sum_{e\in E(A,B)}x_e^2+\sum_{\substack{e,e'\in E(A,B)\\ e\neq e'}}2\,\varepsilon_{A,B,e}\varepsilon_{A,B,e'}x_e x_{e'}.
\]

\begin{lem}\label{lem:length}
Let $G$ be connected and $x\in C_1(G)$, with $x=x_C+x_F$ as in
Lemma~\ref{lem:primitive}. Then
\[
|x_C|^2=\frac1\tau\sum_{\{A,B\}}\cF_{\{A,B\}}(x),
\]
the sum over all two-component spanning forests $\{A,B\}$, where $\tau$ is the
number of spanning trees of $G$.
\end{lem}

\begin{proof}
A *connected spanning unicyclic subgraph* $R$ is a connected spanning subgraph
containing exactly one cycle $C_R$. Deleting any edge of $C_R$ yields a spanning
tree. There is a circulation vector $\omega_R=\sum_{e\in C_R}\delta_{R,e}e$,
$\delta_{R,e}\in\{\pm1\}$, unique up to sign, with $\partial\omega_R=0$; thus
$\omega_R\in\Flow(G)$. Define $\cC_R(x)=|\omega_R\cdot x|^2$, which equals
\[
\cC_R(x)=\sum_{e\in C_R}x_e^2+\sum_{\substack{e,e'\in C_R\\ e\neq e'}}2\,\delta_{R,e}\delta_{R,e'}x_e x_{e'}.
\]
We use three bijections (each is the classical deletion correspondence):
\begin{itemize}
\item For a fixed edge $e$: two-component spanning forests $\{A,B\}$ with
$e\in E(A,B)$ $\leftrightarrow$ spanning trees containing $e$ (add/remove $e$).
\item For a fixed edge $e$: connected spanning unicyclic $R$ with $e\in C_R$
$\leftrightarrow$ spanning trees not containing $e$ (the tree plus $e$ creates a
unique cycle through $e$).
\item For distinct edges $e,e'$: forests $\{A,B\}$ with $e,e'\in E(A,B)$
$\leftrightarrow$ unicyclic $R$ with $e,e'\in C_R$ (delete $e,e'$ from $R$).
Under this bijection one verifies
$\varepsilon_{A,B,e}\varepsilon_{A,B,e'}=-\delta_{R,e}\delta_{R,e'}$:
removing the two edges $e,e'$ of the cycle $C_R$ splits $G$ into the two trees
$A,B$, and the cycle traverses $A,B$ alternately so that consistent orientation
along the cycle ($\delta$) corresponds to opposite cut-orientations
($\varepsilon$).
\end{itemize}
Now sum $\cF_{\{A,B\}}(x)+\cC_R(x)$ over all forests and unicyclics. Each square
term $x_e^2$ appears once for each spanning forest with $e\in E(A,B)$ and once
for each unicyclic with $e\in C_R$, i.e.\ $\tau_e+(\tau-\tau_e)=\tau$ times,
where $\tau_e$ is the number of spanning trees through $e$. Each cross term
$2x_ex_{e'}$ appears with coefficient
$\varepsilon_{A,B,e}\varepsilon_{A,B,e'}$ over forests and
$\delta_{R,e}\delta_{R,e'}$ over unicyclics; by the third bijection these occur
on the same index set with opposite signs, so they cancel in the total. Hence
\begin{equation}\label{eq:fulllen}
\sum_{\{A,B\}}\cF_{\{A,B\}}(x)+\sum_R\cC_R(x)=\tau\sum_e x_e^2=\tau|x|^2.
\end{equation}
Apply \eqref{eq:fulllen} to $x_C$. Since $x_C\in\Cut(G)\otimes\bQ$ is orthogonal
to $\Flow(G)$ and each $\omega_R\in\Flow(G)$, we get $\cC_R(x_C)=0$ for all $R$,
so $\tau|x_C|^2=\sum_{\{A,B\}}\cF_{\{A,B\}}(x_C)$. Finally, $\chi_{A,B}\in
\Cut(G)$, while $x_F\in\Flow(G)\otimes\bQ$ is orthogonal to $\Cut(G)$, so
$\chi_{A,B}\cdot x=\chi_{A,B}\cdot x_C$, giving
$\cF_{\{A,B\}}(x)=\cF_{\{A,B\}}(x_C)$. Combining yields the formula.
\end{proof}

\subsection*{Apex trees}

% Given a weighted tree, the vertex set is a basis for L(T); let V^* denote the
% dual basis of L(T)^*, with v^* dual to v.

Let $T$ be a weighted tree; $V$ is a basis of $L(T)$ and $V^{*}$ the dual basis
of $L(T)^{*}$, with $v^{*}$ dual to $v$.

\begin{lem}\label{lem:apex}
Let $T$ be a weighted tree with no bad vertices, and let $v$ be a vertex with
$w(v)>d(v)$. If $a\in L(T)^{*}$ satisfies $a\cdot v\neq0$, then
$|a|^2\geq|v^{*}|^2$.
\end{lem}

\begin{proof}
Form the *apex tree* $G$ by adjoining to $T$ one new vertex $v_\infty$ joined to
each $u\in V$ by $w(u)-d(u)\geq0$ parallel edges. By \eqref{eq:cob} and the
defining form, the map $u\mapsto\partial^{*}u$ ($u\in V$) is an isometry
$L(T)\xrightarrow{\ \cong\ }\Cut(G)$: indeed for $u\in V$ the degree of $u$ in
$G$ is $d_T(u)+(w(u)-d_T(u))=w(u)$, and for distinct $u,u'\in V$ the multiplicity
$e(u,u')$ equals $1$ if $(u,u')\in E$ and $0$ otherwise, matching the form on
$L(T)$. (Here we use that no vertex is bad, so all multiplicities $w(u)-d(u)$ are
$\geq0$; and $G$ is connected since $T$ is.) Henceforth identify
$L(T)\cong\Cut(G)$ and hence $L(T)^{*}\cong\Cut(G)^{*}$.

By hypothesis $w(v)-d(v)\geq1$, so there is at least one edge $e$ of $G$ joining
$v$ to $v_\infty$. Consider the basis vector $e\in C_1(G)$ and its projection
$e_C\in\Cut(G)^{*}$. We claim $e_C=v^{*}$ under our identification. Indeed for
any $u\in V$, $e_C\cdot\partial^{*}u=e\cdot\partial^{*}u=\partial e\cdot u$, and
$\partial e=v-v_\infty$ (orienting $e$ from $v_\infty$ to $v$), so
$e\cdot\partial^{*}u=v\cdot u-v_\infty\cdot u=\delta_{u,v}$ since $v_\infty\notin
V$. Thus $e_C$ is the dual of $v$, i.e.\ $e_C=v^{*}$.

For a two-component spanning forest $\{A,B\}$ of $G$, the square-flux of $e$ is
$\cF_{\{A,B\}}(e)=|\chi_{A,B}\cdot e|^2$, which equals $1$ if $e\in E(A,B)$ and
$0$ otherwise. Now $e\in E(A,B)$ means the two endpoints of $e$ lie in distinct
trees $A,B$; since $A,B$ are spanning trees of $G$ on complementary vertex sets,
$e\in E(A,B)$ if and only if $A\cup B\cup\{e\}$ is a spanning tree of $G$
containing $e$, and this correspondence is a bijection (delete $e$ from such a
tree). Hence by Lemma~\ref{lem:length},
\begin{equation}\label{eq:vstar}
|v^{*}|^2=|e_C|^2=\frac1\tau\sum_{\{A,B\}}\cF_{\{A,B\}}(e)=\frac{\tau_e}{\tau},
\end{equation}
where $\tau_e$ is the number of spanning trees of $G$ containing $e$.

Now take $a\in L(T)^{*}\cong\Cut(G)^{*}$ with $a\cdot v\neq0$; write $x_C\in
\Cut(G)^{*}$ for its image and choose (by surjectivity of orthogonal
projection, Lemma~\ref{lem:primitive}) some $x\in C_1(G)$ with $x_C$ as cut
component, $x=x_C+x_F$. By Lemma~\ref{lem:length},
$|a|^2=|x_C|^2=\frac1\tau\sum_{\{A,B\}}\cF_{\{A,B\}}(x)$. To prove
$|a|^2\geq|v^{*}|^2=\tau_e/\tau$ it suffices to exhibit at least $\tau_e$
distinct two-component spanning forests $\{A,B\}$ with $\cF_{\{A,B\}}(x)\neq0$;
we do so via an injection $\phi$ from $\{$spanning trees of $G$ containing
$e\}$ into $\{$forests with nonzero square-flux$\}$.

Let $R$ be a spanning tree of $G$ with $e\in R$. Let $R_1,\dots,R_k$ be the
connected components of $R-v$ (here $v$ is the chosen vertex of $T$, and $R-v$
deletes $v$ and its incident $R$-edges). Each component $R_j$, together with its
complementary subtree $R-R_j$ (which contains $v$), forms a two-component
spanning forest $\{R_j,R-R_j\}$; these are exactly the forests obtained by
deleting from $R$ a single edge incident to $v$. Since every vertex of $G-v$
lies in exactly one $R_j$, \eqref{eq:chi} gives the relation in $\Cut(G)$
\[
\chi_{V(G)\setminus\{v\},\,\{v\}}=\sum_{j=1}^{k}\chi_{R_j,\,R-R_j},
\]
because $-\partial^{*}v=\sum_{u\in V(G)\setminus\{v\}}(\dots)$ matches the
sum of the cut vectors over the components (both sides equal
$-\sum_{u\neq v}\partial^{*}u = \partial^* v$ in $\Cut(G)$; more directly each
side is the cut vector of the partition separating $v$ from everything else,
decomposed over the components $R_j$). Pairing with $x$,
\[
\chi_{V(G)\setminus\{v\},\{v\}}\cdot x=\sum_{j}\chi_{R_j,R-R_j}\cdot x .
\]
The left side equals $-\partial^{*}v\cdot x=-x_C\cdot\partial^{*}v=-a\cdot
v\neq0$. Hence some index $\ell$ has $\chi_{R_\ell,R-R_\ell}\cdot x\neq0$, i.e.\
$\cF_{\{R_\ell,R-R_\ell\}}(x)\neq0$. Set $\phi(R)=\{R_\ell,R-R_\ell\}$ for one
such $\ell$.

We check $\phi$ is injective by reconstructing $R$ from $\phi(R)=\{A,B\}$. By
construction $\{A,B\}$ is obtained from $R$ (a tree containing $e$) by deleting a
single edge incident to $v$. If $e\in E(A,B)$, then $A,B$ are precisely the two
components after deleting $e$, so $R=A\cup B\cup\{e\}$. Otherwise $e$ lies inside
one part, say $A$; then $v$ and $v_\infty$ both lie in $A$, so $B$ is a subtree
of $T$ not containing $v$. Since $T$ is a tree, there is a unique edge $f$ of $T$
joining $v$ to a vertex of $B$, and $R=A\cup B\cup\{f\}$. In either case $R$ is
determined by $\{A,B\}$, so $\phi$ is injective. (Distinct trees in the same
fiber would force the same reconstruction, contradicting distinctness.)

Therefore there are at least $\tau_e$ forests with nonzero square-flux, giving
$|a|^2\geq\tau_e/\tau=|v^{*}|^2$.
\end{proof}

\begin{cor}\label{cor:apex}
Under the hypotheses of Lemma~\ref{lem:apex}, $v^{*}\in L(T)^{*}$ is irreducible
in $L(T)^{*}$.
\end{cor}

\begin{proof}
Suppose $v^{*}=a+b$ with $a,b\in L(T)^{*}$ nonzero. Then
$1=v^{*}\cdot v=a\cdot v+b\cdot v$, so without loss of generality $a\cdot
v\neq0$. By Lemma~\ref{lem:apex}, $|a|^2\geq|v^{*}|^2$, and $|b|^2>0$ since
$b\neq0$ and the form is positive definite (note $L(T)^{*}\otimes\bQ=L(T)\otimes
\bQ$ carries the positive definite form). Then
\[
|v^{*}|^2=|a+b|^2=|a|^2+2(a\cdot b)+|b|^2>|v^{*}|^2+2(a\cdot b),
\]
forcing $a\cdot b<0$. Hence no decomposition with $a\cdot b\geq0$ exists, i.e.\
$v^{*}$ is irreducible.
\end{proof}

\begin{cor}\label{cor:tight}
Under the hypotheses of Lemma~\ref{lem:apex}, $0<|v^{*}|^2\leq1$, and if
$|v^{*}|^2=1$ then either $w(v)=1$ or $T$ has a leaf $u\neq v$ with $w(u)=1$.
\end{cor}

\begin{proof}
From \eqref{eq:vstar}, $|v^{*}|^2=\tau_e/\tau$ with $1\le\tau_e\le\tau$, so
$0<|v^{*}|^2\leq1$. Now $|v^{*}|^2=1\iff\tau_e=\tau\iff$ every spanning tree of
$G$ uses $e\iff e$ is a cut-edge (bridge) of $G$. A bridge cannot lie on any
cycle; since $e$ joins $v$ to the apex $v_\infty$, this happens precisely when
$e$ is the unique edge of $G$ incident to $v_\infty$, i.e.\ $w(v)-d(v)=1$ and
$w(u)-d(u)=0$ for all $u\neq v$. (If $v_\infty$ had another incident edge $e'$,
then together with paths through $T$ one finds a cycle through $e$, since
$G-\{e\}$ would still connect $v$ to $v_\infty$ via $e'$ and $T$, contradicting
that $e$ is a bridge.) Thus $w(u)=d(u)$ for every $u\neq v$. If $V=\{v\}$ then
$d(v)=0$ and $w(v)=1$. Otherwise $T$ has a leaf $u\neq v$, and then
$w(u)=d(u)=1$.
\end{proof}

\subsection*{Reducible vertices}

% Now suppose L(T) is positive definite. Fix a vertex v and analyze when v is
% reducible by passing to the subtrees obtained by deleting v.

Assume $L(T)$ is positive definite. Fix $v\in V$ and let $T_1,\dots,T_k$ be the
weighted trees forming the connected components of $T-v$ (with weights inherited
from $T$, so $w_{T_i}=w_T|_{V(T_i)}$). For each $i$ let $v_i\in T_i$ be the
unique vertex adjacent to $v$ in $T$. Let $\lambda_i\in L(T_i)^{*}$ be the dual
of $v_i$ with respect to the basis $V(T_i)$. Via the chain of inclusions
$L(T_i)^{*}\subset L(T_i)\otimes\bQ\subset L(T)\otimes\bQ$ we view $\lambda_i\in
L(T)\otimes\bQ$.

Set $L_1=L(T_1)\oplus\cdots\oplus L(T_k)\subseteq L(T)$, spanned by
$V\setminus\{v\}$; since this basis completes to $V$, $L_1$ is a primitive
sublattice. Its orthogonal complement $L_1^{\perp}\subset L(T)\otimes\bQ$ is
one-dimensional. Let $\pi:L(T)\to L_1^{\perp}$ be orthogonal projection and
$L_0=\pi(L(T))$, a rank-one rational lattice. Then
$L(T)\subseteq L_0\oplus L_1^{*}=L_0\oplus L(T_1)^{*}\oplus\cdots\oplus
L(T_k)^{*}$, and every $z\in L(T)$ decomposes orthogonally as
$z=\pi(z)+\sum_i z_i$ with $z_i\in L(T_i)^{*}$ the projection of $z$ to
$L(T_i)\otimes\bQ$.

\begin{lem}\label{lem:irred0}
$\pi(v)$ generates $L_0$; in particular $\pi(v)$ is irreducible in $L_0$.
\end{lem}

\begin{proof}
Every $u\in V\setminus\{v\}$ lies in $L_1$, so $\pi(u)=0$. Since $V$ generates
$L(T)$, the image $L_0=\pi(L(T))$ is generated by $\pi(v)$. A generator of a
rank-one positive definite lattice is irreducible: if $\pi(v)=a+b$ with
$a,b\in L_0$ nonzero, then $a=ma_0$, $b=na_0$ for the generator $a_0$ with
$m+n=\pm1$ and $m,n\neq0$, whence $a\cdot b=mn|a_0|^2<0$ (as $mn<0$). So $\pi(v)$
is irreducible.
\end{proof}

\begin{lem}\label{lem:dichotomy}
If $v$ is reducible in $L(T)$, then there is an index $i$ such that either
\textup{(i)} $\lambda_i$ is reducible in $L(T_i)^{*}$, or \textup{(ii)}
$\lambda_i\in L(T_i)$.
\end{lem}

\begin{proof}
For each $i$ and each $u\in V(T_i)$ we have $v\cdot u=-1$ if $u=v_i$ and
$v\cdot u=0$ otherwise; thus the projection of $v$ onto $L(T_i)\otimes\bQ$ is the
unique element $\mu\in L(T_i)^{*}$ with $\mu\cdot u=v\cdot u$ for all $u\in
V(T_i)$, namely $\mu=-v_i^{*,T_i}=-\lambda_i$. Hence the orthogonal
decomposition of $v$ is
\begin{equation}\label{eq:vdecomp}
v=\pi(v)-\lambda_1-\cdots-\lambda_k.
\end{equation}

Suppose $v=x+y$ with nonzero $x,y\in L(T)$ and $x\cdot y\geq0$. Decompose
$x=\bar x+\sum_i x_i$ and $y=\bar y+\sum_i y_i$ with $\bar x,\bar y\in L_0$ and
$x_i,y_i\in L(T_i)^{*}$. Then $\bar x+\bar y=\pi(v)$ and $x_i+y_i=-\lambda_i$ for
all $i$. By orthogonality of the summands,
\[
0\leq x\cdot y=\bar x\cdot\bar y+\sum_{i=1}^{k}x_i\cdot y_i.
\]
By Lemma~\ref{lem:irred0}, $\pi(v)$ is irreducible in $L_0$; since
$\bar x+\bar y=\pi(v)$, either one of $\bar x,\bar y$ is $0$ (then $\bar x\cdot
\bar y=0$) or both are nonzero, in which case irreducibility forces $\bar x\cdot
\bar y<0$. In all cases $\bar x\cdot\bar y\le0$, with equality only if
$\bar x=0$ or $\bar y=0$.

\emph{Case 1: $\bar x\cdot\bar y<0$.} Then $\sum_i x_i\cdot y_i\geq-\bar x\cdot
\bar y>0$, so some $x_i\cdot y_i>0$. Since $x_i+y_i=-\lambda_i$, the element
$-\lambda_i$ (equivalently $\lambda_i$, as reducibility is sign-invariant: if
$-\lambda_i=a+b$ with $a\cdot b\ge0$ then $\lambda_i=(-a)+(-b)$ with
$(-a)\cdot(-b)\ge0$) is reducible in $L(T_i)^{*}$, provided $x_i,y_i$ are both
nonzero; and $x_i\cdot y_i>0$ forces both nonzero. This gives conclusion (i).

\emph{Case 2: $\bar x\cdot\bar y=0$ and all $x_i\cdot y_i\le0$.} If some
$x_i\cdot y_i>0$ we are in Case~1's conclusion (i), so suppose every $x_i\cdot
y_i\le0$. Then $0\le x\cdot y=\bar x\cdot\bar y+\sum_i x_i\cdot y_i\le0$ forces
$\bar x\cdot\bar y=0$ and $x_i\cdot y_i=0$ for all $i$. If some $x_i\cdot y_i>0$
were possible we'd be in (i); otherwise, from $\bar x\cdot\bar y=0$ and the
irreducibility of $\pi(v)$ we get $\{\bar x,\bar y\}=\{0,\pi(v)\}$, and from
$x_i\cdot y_i=0$ with $x_i+y_i=-\lambda_i$: if both $x_i,y_i$ were nonzero we
would have a decomposition of $-\lambda_i$ with $x_i\cdot y_i=0\geq0$, i.e.\
$\lambda_i$ reducible, giving (i); so assume $\{x_i,y_i\}=\{0,-\lambda_i\}$ for
each $i$.

Thus, after possibly relabeling so that $\bar x=0$, we have
$x=-\sum_{i\in A}\lambda_i$ for some $A\subseteq\{1,\dots,k\}$, and $A\neq
\emptyset$ since $x\neq0$. Now $x\in L(T)$, and $x\in L_1\otimes\bQ$ (as each
$\lambda_i\in L(T_i)\otimes\bQ\subseteq L_1\otimes\bQ$). Since $L_1$ is primitive
in $L(T)$, $L(T)\cap(L_1\otimes\bQ)=L_1$, so $x\in L_1=\bigoplus_i L(T_i)$. The
summands $L(T_i)$ are orthogonal, and the component of $x$ in $L(T_i)\otimes\bQ$
is $-\lambda_i$ for $i\in A$ and $0$ otherwise. As $x\in\bigoplus_i L(T_i)$, its
$L(T_i)$-component lies in $L(T_i)$; hence $\lambda_i\in L(T_i)$ for every $i\in
A\neq\emptyset$. This is conclusion (ii).
\end{proof}

\subsection*{Proof of the Theorem}

\begin{proof}[Proof of Theorem~\ref{thm:main}]
We prove the stronger statement: either the unique bad vertex $v$ is
irreducible in $L(T)$, or $T$ contains a leaf of weight $1$ (which is
automatically irreducible, being a vertex element $u$ with $|u|^2=w(u)=1$, hence
of minimal nonzero norm; explicitly if $u=a+b$ with $a,b\neq0$ then
$1=|u|^2=|a|^2+2a\cdot b+|b|^2\geq 1+1+2a\cdot b>1+2a\cdot b$ when
$|a|^2,|b|^2\ge1$, but actually $|a|^2,|b|^2\geq1$ as integral positive definite
norms of nonzero vectors, forcing $a\cdot b\le -\tfrac12<0$).

Let $v$ be the unique bad vertex, and define $T_1,\dots,T_k$, $v_i$,
$\lambda_i$ as above. Since $v$ is the only bad vertex of $T$, each $T_i$
contains no bad vertex (the deletion of $v$ does not increase any weight nor
decrease degrees of vertices other than $v_i$, and only $v_i$ loses a unit of
degree, which only helps the inequality $w\geq d$; vertices in $T_i$ other than
$v_i$ retain their $T$-degree and weight and were good in $T$). Moreover, in
$T_i$ the degree of $v_i$ is one less than its degree $d_T(v_i)$ in $T$, and
$d_T(v_i)\leq w(v_i)$ because $v_i\neq v$ is good in $T$. Hence
\[
d_{T_i}(v_i)=d_T(v_i)-1\leq w(v_i)-1<w(v_i),
\]
so $w(v_i)>d_{T_i}(v_i)$. Therefore $T_i$ and the vertex $v_i$ satisfy the
hypotheses of Lemma~\ref{lem:apex} and Corollaries~\ref{cor:apex},
\ref{cor:tight} (note also that $L(T_i)$ is positive definite, as a sublattice
of the positive definite $L(T)$).

By Corollary~\ref{cor:apex}, $\lambda_i$ is irreducible in $L(T_i)^{*}$ for
every $i$.

Now suppose first that $\lambda_i\notin L(T_i)$ for all $i$. If $v$ were
reducible, Lemma~\ref{lem:dichotomy} would supply an index $i$ with either
$\lambda_i$ reducible in $L(T_i)^{*}$ (contradicting
Corollary~\ref{cor:apex}) or $\lambda_i\in L(T_i)$ (contrary to our standing
assumption). Hence $v$ is irreducible, and the theorem holds with $u=v$.

Otherwise $\lambda_i\in L(T_i)$ for some $i$. Then $|\lambda_i|^2\in\bZ$, while
$|\lambda_i|^2=|v_i^{*,T_i}|^2$ satisfies $0<|\lambda_i|^2\leq1$ by
Corollary~\ref{cor:tight}; therefore $|\lambda_i|^2=1$. Corollary~\ref{cor:tight}
then asserts that either $w(v_i)=1$ or $T_i$ contains a leaf $u\neq v_i$ with
$w(u)=1$. In the first case $v_i$ is a vertex of $T$ with weight $1$; since
$d_{T_i}(v_i)=0$ when $w(v_i)=1=|\lambda_i|^2$ forces (by the proof of
Corollary~\ref{cor:tight}) the edge to be a bridge, in fact $v_i$ is a leaf of
$T$ of weight $1$ when $T_i=\{v_i\}$, and in general $w(v_i)=1$ with
$d_{T_i}(v_i)\geq 0$ together with $|\lambda_i|^2=1$ implies $v_i$ is a leaf of
$T_i$, hence has $T$-degree $d_{T_i}(v_i)+1$; but $w(v_i)=1$ and $v_i$ good in
$T$ gives $d_T(v_i)\le1$, so $d_{T_i}(v_i)\le0$, i.e.\ $T_i=\{v_i\}$ and $v_i$ is
a leaf of $T$ of weight $1$. In the second case, $u$ is a leaf of $T_i$, and
since $u\neq v_i$ it is not adjacent to $v$ in $T$, so $u$ is a leaf of $T$ as
well, of weight $1$. In either case $T$ contains a leaf of weight $1$, which is
irreducible as noted at the outset.

In all cases $T$ contains an irreducible vertex element, completing the proof.
\end{proof}

% Remark: the example T with central vertex of weight 2, degree 3, and leaves of
% weights 1,2,3 shows the bad vertex itself can be reducible, so the dichotomy
% (and the appeal to a weight-1 leaf) is genuinely necessary.